\scanpage{79}
\[
\alpha:\underline{\omega}_J\simeq\mu_4^*
\qquad\text{i.e. d'un élément de }\underline{\omega}_J\wedge\mu_4^*\,.
\tag*{(20)}
\]

Ces définitions posées, on voit donc que la catégorie\inserted{-groupoïde} des courbes elliptiques\inserted{$/k$} à \unsure{épinglages|épinglage} de Legendre \deleted{\unsure{avec|est}} \inserted{est} équivalente à la catégorie-groupoïde des \deleted{triplets couples} \inserted{triplets} $(J,t,\alpha)$ \deleted{de} formés

\[
\left\{\begin{aligned}
&\text{a) d'un ens }J\text{ à trois éléments (d'où }\Sigma_J,D_J,\Sigma_J^*,\\
&\text{b) d'un } t\in\Sigma_J^*(k)\\
&\text{c) d'un isom }\alpha:\underline{\omega}_J\simeq\mu_4^{\unsure{*|'}}\deleted{\unsure{(k)}}\quad(=\mathord{\mathrm{I}\mkern-7mu\mu}_4^*(k))
\end{aligned}\right.
\tag*{(21)}
\]

Un épinglage de Legendre-Jacobi d'échelon 2 \inserted{d'une courbe elliptique $E$} sera la donnée d'un couple d'épinglages, l'un de Legendre, l'autre de Jacobi d'échelon 2 i.e. d'une bijection $J\simeq\{0,1,\infty\}$, permettant d'identifier $\Sigma_J$ à la droite projective type $\mathbb{P}^1_{\unsure{\mathbb{C}|k}}$. On a alors dans $\underline{\omega}_J\simeq\underline{\omega}_{\{0,1,\infty\}}$ un élément privilégié $\omega^+$, et la donnée c) \deleted{\unsure{du dessus|du choix est|du dessin}} de l'épinglage de Legendre équivaut \inserted{donc} à celle d'une racine carrée \deleted{\unsure{de l'unité|de $-1$|unité}} de $-1$, soit $i$ -- c'est la donnée ``\unsure{arithmétique|authentique}''. \deleted{Par ailleurs, la catég} La catégorie-groupoïde des courbes elliptiques sur $k$, munies d'un épinglage de Legendre-Jacobi, \deleted{équivaut à celle des \unsure{ses|ces} objets à isom. près ca} \inserted{est rigide (pas d'auto. non-identiques)} et elle
% archive sheet number: 78
